\section{Introduction}
\label{sec:introduction}

Passwordless authentication built on FIDO2, comprising the WebAuthn API~\cite{webauthn-l3} and the Client to Authenticator Protocol (CTAP)~\cite{ctap21}, is now a default authentication option across major browsers, operating systems, and online services~\cite{fido2-overview}. The authenticator side of this protocol, the device that holds credentials and proves possession of them, has almost universally been implemented as a dedicated security key: a firmware image whose sole purpose is to be an authenticator.

That assumption increasingly does not match how embedded deployments actually look. A growing class of MCU-based devices, sensor nodes, industrial controllers, smart-lock firmware, and similar shared-purpose systems, would benefit from FIDO2 authentication as one feature among several, not as the reason the firmware exists. Existing open FIDO2 implementations such as OpenSK~\cite{opensk} and SoloKeys/Solo2~\cite{solokeys} are built around the dedicated-device assumption and do not address this deployment shape. Retrofitting a security-critical, externally-specified protocol into firmware that was not designed around it, without requiring the host application to know or care about the authenticator subsystem's internals, and without leaking security-critical state across that boundary, is an under-studied integration scenario.

Prior work addresses adjacent parts of this problem but not its intersection. Protocol-security analyses establish that FIDO2 is cryptographically sound as specified, independent of how any particular implementation is structured as software. Hardware-based MCU compartmentalization techniques isolate components using memory protection or trusted execution, a different mechanism than software boundary isolation and one that assumes hardware most MCUs in this deployment class do not have. Zephyr and RTOS subsystem case studies demonstrate that a real subsystem contribution can itself be academic evidence, but none of them operate in a security-critical, externally-specified protocol domain, and none pair architectural claims with a falsifiability apparatus that ties those claims to structurally checkable evidence in the source tree. No existing work sits at the intersection of a software-boundary-isolated, falsifiably-evidenced, security-critical subsystem retrofitted into shared-purpose firmware. That intersection is the gap this paper addresses.

This paper's contribution is a reusable architectural pattern: strict interface boundaries between transport, crypto, CBOR parsing, storage, attestation, and user-presence concerns, such that each is independently replaceable and none has implicit knowledge of the others' internals. Security here is a property of the subsystem's boundary, not a claim about cryptographic primitives: the underlying cryptographic operations are delegated to the PSA Crypto API~\cite{psa-crypto-api} and are out of scope, while what the subsystem itself owns and must be shown correct is the isolation logic around that boundary, retry-counter ordering, lockout state, credential binding parameters, and the guarantee that a buggy or compromised host application cannot reach that state directly. The pattern is demonstrated, not merely proposed, through a complete CTAP2.1 implementation built as a Zephyr RTOS subsystem. The implementation is the vehicle; the pattern, and the evidence that it holds under measurement rather than only under inspection, is the contribution.

Three research questions structure the evaluation.

\begin{itemize}
\item \textbf{RQ1.} Composability: can the subsystem be integrated as a side feature into an existing, unrelated Zephyr application without touching that application's core logic?
\item \textbf{RQ2.} Portability: does the boundary design deliver real hardware portability, not just architectural intent? A second transport, a Bluetooth Low Energy implementation authored independently against the existing transport API, is used as the primary evidence: does it require zero modifications to the protocol core?
\item \textbf{RQ3.} Cost of isolation: what is the measurable overhead, in footprint, latency, and dispatch cost, of enforcing these boundaries compared to tightly coupled alternatives, including a comparison against SoloKeys and, conditionally, OpenSK?
\end{itemize}

The generalizability claim in this paper is scoped deliberately. RQ1 and RQ2 jointly support the claim that this architectural pattern works for FIDO2 as an OS-level subsystem, not a one-off application. They do not claim that the pattern generalizes to a second, unrelated subsystem; no second prototype is built or evaluated here, and that extension is left as future work.

The remainder of this paper is organized as follows. Section~\ref{sec:related} reviews related work. Section~\ref{sec:architecture} presents the system architecture. Section~\ref{sec:design-principles} details the design principles underlying the boundary structure. Section~\ref{sec:security} discusses security considerations. Section~\ref{sec:concurrency} describes the concurrency and real-time behavior of the subsystem. Section~\ref{sec:methodology} states the research questions and evaluation methodology. Section~\ref{sec:evaluation} presents evaluation results. Section~\ref{sec:discussion} discusses what the results support. Section~\ref{sec:limitations} states limitations and scope. Section~\ref{sec:conclusion} concludes.